% GENERATED FILE. Edit canonical lessons, then run npm run generate:books.
% canonical-source-hash: fnv1a64:9fb4e2d7d230f85b
% canonical-lessons: IT-C369-partito, IT-C369-parlamento, IT-C369-ministro, IT-C369-eleggere, IT-C369-elettore

\chapter{La politica}
\label{ch:it-la-politica}
\hlchaptermodality{369}

\begin{chapteropening}
\textbf{By the end of this chapter:} \emph{I can talk about politics: a party, parliament, a minister, voters and electing someone.}

Complete the last lesson of chapter 369: I can talk about politics: a party, parliament, a minister, voters and electing someone.
\end{chapteropening}

\section[il partito]{il partito — the political party}

\label{lesson:IT-C369-partito}



\begin{quote}
Before the new one: say the Italian for to respect, then the Italian for the integration.
\end{quote}

\subsection*{You'll want to know: il partito}
\textbf{il partito} — \textquotedblleft{}the political party\textquotedblright{}.

\begin{quote}
Non sono iscritto a nessun partito.
\end{quote}

\begin{quote}
\emph{I am not a member of any party.}
\end{quote}

Also the past participle of \emph{partire}: \emph{è partito}.

\begin{cousinweb}
\emph{partito} is the past participle of \emph{partire} in its old sense \textquotedblleft{}to divide\textquotedblright{}, inherited from Latin \emph{partīrī} \textquotedblleft{}to divide\textquotedblright{}, from \emph{pars, partis} \textquotedblleft{}a part\textquotedblright{} (Lewis \& Short): one part of the people.
\end{cousinweb}

\begin{grammarlens}[title={What you've built}]
One more word for telling, explaining and describing at B1.
\end{grammarlens}

\subsection*{Your turn}
\begin{itemize}
\raggedright
  \item \emph{Say it:} \emph{il partito}
  \item \emph{Say it:} \emph{il partito}, once more
  \item \emph{Say it:} a sentence of your own with \emph{il partito}
  \item \emph{Recall:} say the Italian for to respect, then the Italian for the integration, then say \emph{il partito} again
\end{itemize}

\begin{tcolorbox}[breakable,colback=teal!4,colframe=teal!35!black,title={Before you move on}]
What does il partito mean? (\textquotedblleft{}The political party\textquotedblright{}.) Say it once more.
\end{tcolorbox}
\section[il parlamento]{il parlamento — the parliament}

\label{lesson:IT-C369-parlamento}



\begin{quote}
Before the new one: say the Italian for the integration, then the Italian for the political party.
\end{quote}

\subsection*{You'll want to know: il parlamento}
\textbf{il parlamento} — \textquotedblleft{}the parliament\textquotedblright{}.

\begin{quote}
Il parlamento discute la nuova legge.
\end{quote}

\begin{quote}
\emph{Parliament is debating the new law.}
\end{quote}

Stress: \emph{parlaménto}. Its members are elected by the citizens.

\begin{cousinweb}
\emph{parlamento} is built on \emph{parlare} \textquotedblleft{}to speak\textquotedblright{}, inherited from Late Latin \emph{parabolāre} (EtymDB), after Old French \emph{parlement} \textquotedblleft{}a speaking, a meeting\textquotedblright{}: the place where people talk.
\end{cousinweb}

\begin{grammarlens}[title={What you've built}]
One more word for telling, explaining and describing at B1.
\end{grammarlens}

\subsection*{Your turn}
\begin{itemize}
\raggedright
  \item \emph{Say it:} \emph{il parlamento}
  \item \emph{Say it:} \emph{il parlamento}, once more
  \item \emph{Say it:} a sentence of your own with \emph{il parlamento}
  \item \emph{Recall:} say the Italian for the integration, then the Italian for the political party, then say \emph{il parlamento} again
\end{itemize}

\begin{tcolorbox}[breakable,colback=teal!4,colframe=teal!35!black,title={Before you move on}]
What does il parlamento mean? (\textquotedblleft{}The parliament\textquotedblright{}.) Say it once more.
\end{tcolorbox}
\section[il ministro]{il ministro — the minister}

\label{lesson:IT-C369-ministro}



\begin{quote}
Before the new one: say the Italian for the political party, then the Italian for the parliament.
\end{quote}

\subsection*{You'll want to know: il ministro}
\textbf{il ministro} — \textquotedblleft{}the minister\textquotedblright{}.

\begin{quote}
Il ministro della salute parla in televisione stasera.
\end{quote}

\begin{quote}
\emph{The health minister is speaking on television tonight.}
\end{quote}

\emph{Il primo ministro} — the prime minister. A woman is also \emph{la ministra}.

\begin{cousinweb}
\emph{ministro} is a learned borrowing from Latin \emph{minister} \textquotedblleft{}servant, attendant\textquotedblright{} (Lewis \& Short), built on \emph{minus} \textquotedblleft{}less\textquotedblright{}: first the one who serves, then the one who serves the state.
\end{cousinweb}

\begin{grammarlens}[title={What you've built}]
One more word for telling, explaining and describing at B1.
\end{grammarlens}

\subsection*{Your turn}
\begin{itemize}
\raggedright
  \item \emph{Say it:} \emph{il ministro}
  \item \emph{Say it:} \emph{il ministro}, once more
  \item \emph{Say it:} a sentence of your own with \emph{il ministro}
  \item \emph{Recall:} say the Italian for the political party, then the Italian for the parliament, then say \emph{il ministro} again
\end{itemize}

\begin{tcolorbox}[breakable,colback=teal!4,colframe=teal!35!black,title={Before you move on}]
What does il ministro mean? (\textquotedblleft{}The minister\textquotedblright{}.) Say it once more.
\end{tcolorbox}
\section[eleggere]{eleggere — to elect}

\label{lesson:IT-C369-eleggere}



\begin{quote}
Before the new one: say the Italian for the parliament, then the Italian for the minister.
\end{quote}

\subsection*{You'll want to know: eleggere}
\textbf{eleggere} — \textquotedblleft{}to elect\textquotedblright{}.

\begin{quote}
I cittadini eleggono il sindaco ogni cinque anni.
\end{quote}

\begin{quote}
\emph{The citizens elect the mayor every five years.}
\end{quote}

Conjugated like \emph{leggere}; past participle \emph{eletto}.

\begin{cousinweb}
\emph{eleggere} comes from Latin \emph{ēligere} \textquotedblleft{}to pick out, to choose\textquotedblright{} (EtymDB), built from \emph{ē} \textquotedblleft{}out\textquotedblright{} and \emph{legere} \textquotedblleft{}to gather, to choose\textquotedblright{}.
\end{cousinweb}

\begin{grammarlens}[title={What you've built}]
One more word for telling, explaining and describing at B1.
\end{grammarlens}

\subsection*{Your turn}
\begin{itemize}
\raggedright
  \item \emph{Say it:} \emph{eleggere}
  \item \emph{Say it:} \emph{eleggere}, once more
  \item \emph{Say it:} a sentence of your own with \emph{eleggere}
  \item \emph{Recall:} say the Italian for the parliament, then the Italian for the minister, then say \emph{eleggere} again
\end{itemize}

\begin{tcolorbox}[breakable,colback=teal!4,colframe=teal!35!black,title={Before you move on}]
What does eleggere mean? (\textquotedblleft{}To elect\textquotedblright{}.) Say it once more.
\end{tcolorbox}
\section[l'elettore]{l'elettore — the voter}

\label{lesson:IT-C369-elettore}



\begin{quote}
Before the new one: say the Italian for the minister, then the Italian for to elect.
\end{quote}

\subsection*{You'll want to know: l'elettore}
\textbf{l'elettore} — \textquotedblleft{}the voter\textquotedblright{}.

\begin{quote}
Molti elettori non sono andati a votare.
\end{quote}

\begin{quote}
\emph{Many voters did not go and vote.}
\end{quote}

A woman is \emph{l'elettrice}.

\begin{cousinweb}
\emph{elettore} is a learned borrowing from Latin \emph{ēlēctor} \textquotedblleft{}one who chooses\textquotedblright{} (EtymDB), from \emph{ēligere} \textquotedblleft{}to pick out\textquotedblright{}, built on \emph{legere} \textquotedblleft{}to choose\textquotedblright{}.
\end{cousinweb}

\begin{grammarlens}[title={What you've built}]
One more word for telling, explaining and describing at B1.
\end{grammarlens}

\subsection*{Your turn}
\begin{itemize}
\raggedright
  \item \emph{Say it:} \emph{l'elettore}
  \item \emph{Say it:} \emph{l'elettore}, once more
  \item \emph{Say it:} a sentence of your own with \emph{l'elettore}
  \item \emph{Recall:} say the Italian for the minister, then the Italian for to elect, then say \emph{l'elettore} again
\end{itemize}

\begin{tcolorbox}[breakable,colback=teal!4,colframe=teal!35!black,title={Before you move on}]
What does l'elettore mean? (\textquotedblleft{}The voter\textquotedblright{}.) Say it once more.
\end{tcolorbox}
